\begin{abstract}

While multilingual Retrieval-Augmented Generation (mRAG) systems preferentially cite English sources, it remains unclear whether models actually trust English facts more when sources directly contradict each other. This study investigates whether multilingual Large Language Models (LLMs) systematically favor claims written in English over Korean when resolving direct factual contradictions, regardless of query language or prior knowledge. Using a counterbalanced setup across culturally neutral question-answer pairs from the MKQA benchmark, each question is evaluated on conflicting claim pairs and their language-swapped counterparts. By testing model behavior across both English and Korean queries with uniform sentence templates, the framework isolates language preference from factual accuracy to determine whether citation bias extends into actual factual trust.

\end{abstract}